\documentclass[11pt,a4paper]{article}
\usepackage{fontspec}
\setmainfont{TeX Gyre Pagella}
\newfontfamily\headfont{Inter}[UprightFont=*-ExtraBold,BoldFont=*-ExtraBold,Extension=.otf,Path=/usr/share/fonts/opentype/inter/]
\usepackage[a4paper,margin=2.3cm]{geometry}
\usepackage{graphicx,booktabs,float,caption,array,xcolor}
\usepackage[hidelinks]{hyperref}
\usepackage{titlesec}
\usepackage{setspace}
\setstretch{1.05}
\renewcommand{\topfraction}{0.9}
\renewcommand{\bottomfraction}{0.8}
\renewcommand{\textfraction}{0.1}
\renewcommand{\floatpagefraction}{0.75}
\setlength{\parindent}{0pt}
\setlength{\parskip}{0.55em}
\titleformat{\section}{\headfont\Large}{\thesection}{0.8em}{}
\titleformat{\subsection}{\headfont\large}{\thesubsection}{0.8em}{}
\captionsetup{font=small,labelfont=bf}
\urlstyle{same}
\hypersetup{pdftitle={Darkest Skies: Where You Can Still See the Milky Way},pdfauthor={Tristan Hotham}}

\begin{document}
\begin{center}
{\Large\bfseries Darkest Skies: Where You Can Still See the Milky Way\\[2pt] A Sky Glow Model of the World and the United Kingdom from VIIRS Night Lights, Checked Against Citizen Sky Meter Readings\par}
\vspace{1.2em}
{\large Tristan Hotham\par}
\vspace{0.3em}
{\small Independent researcher\\
ORCID: 0000-0003-4931-9349\\
DOI: \href{https://doi.org/10.5281/zenodo.23757708}{https://doi.org/10.5281/zenodo.23757708}\par}
\vspace{0.8em}
2026
\end{center}

\begin{quote}
\small
\begin{center}\textbf{Abstract}\end{center}
Artificial light scattered in the atmosphere brightens the night sky and hides the Milky Way from much of humanity. I build a simple, open sky glow model from the Earth Observation Group's annual VIIRS night lights for 2023, read for all {{meta.world_files|int}} countries from WorldPop's harmonised 3 arc second files, and use it to map artificial zenith brightness for the world at 2.5 arc minutes and for the United Kingdom at 30 arc seconds. Each lit cell's light spreads with a power of distance; the exponent and the scale are fitted to {{model.n18|int}} Globe at Night Sky Quality Meter readings from 2019 to 2025, and in spatial cross-validation the model misses a typical reading by {{model.cv18_mae|f2}} magnitudes per square arcsecond (correlation {{model.cv18_r|f2}}). With WorldPop's 2023 population and the thresholds of the 2016 World Atlas of Artificial Night Sky Brightness, {{world.pop_mw|pct}} of the world's people live where artificial glow hides the Milky Way, against {{falchi.pop_mw|pct}} in the atlas for 2014; the two agree closely on the ranking of {{falchi.compare.n|int}} countries (Spearman's $\rho$ = {{falchi.compare.spearman|f2}}). In the United Kingdom the share is {{uk.pop_mw|pct}} of people but only {{uk.land_mw|pct}} of land. The darkest areas are the Scottish islands and Highlands; the nearest Milky Way sky is {{uk_cities.London.mw_km|int}} km from central London. Radiance seen from space fell {{change.gbr_pct|abs1}}\% over Great Britain between 2015 and 2023, a change that the sensor's blindness to blue LED light probably exaggerates.

\textbf{Keywords:} light pollution; night sky brightness; sky glow; Milky Way; VIIRS; night lights; Globe at Night; Sky Quality Meter; United Kingdom; data visualisation
\end{quote}

\section{Introduction}

Away from towns, on a clear night with no Moon, the Milky Way is the brightest diffuse feature of the sky. It fades quickly once artificial light is added: light leaving streets and buildings is scattered back down by air molecules and aerosols and raises the brightness of the whole sky, so that the faint, extended glow of the Galaxy no longer stands out. The World Atlas of Artificial Night Sky Brightness (Falchi et al. 2016) put numbers on this for the whole planet. Using satellite night lights from 2014, a radiative transfer model of how light propagates in the atmosphere (Cinzano and Falchi 2012) and calibration against sky brightness measurements, it found that more than a third of humanity cannot see the Milky Way from where they live.

Ten years on, the question is worth asking again with open, current data, and at a finer scale for one country. This note does three things. It builds a deliberately simple sky glow model from the 2023 VIIRS annual night lights, so that every step can be rerun from open files; it tests the model against thousands of hand-held sky meter readings collected by volunteers; and it uses the model to answer plain questions for the world and the United Kingdom: who can still see the Milky Way from home, where the darkest places are, and how far city dwellers must go for a dark sky. The results are published with an interactive map at \url{https://maps.tristanhotham.com/dark-skies/}.

\section{Data}

\subsection{Night lights}
The Earth Observation Group (Colorado School of Mines) produces annual composites of the VIIRS Day/Night Band from cloud-free monthly averages, with background noise, gas flares and volcanoes removed (VNL version 2; Elvidge et al. 2021). The sensor passes at about 1:30 am local time and measures radiance in nW/cm\textsuperscript{2}/sr between roughly 500 and 900 nm. WorldPop republishes these composites per country, resampled to 3 arc seconds, as part of its harmonised covariates (Woods et al. 2024), under CC BY 4.0. I read all {{meta.world_files|int}} country files for 2023 in a web browser and reduced them there to a world grid of 2.5 arc minutes (the mean radiance of the valid pixels in each cell and their number), and the files for the United Kingdom, Ireland and the Isle of Man for every year from 2015 to 2023 to a grid of 30 arc seconds.

\subsection{Population}
People are taken from WorldPop's 2023 population grid at 30 arc seconds (Global2 release R2025A, constrained to built-up areas; Bondarenko et al. 2025), summed to the two model grids. The world total is {{world.pop_total_bn|f2}} billion.

\subsection{Sky meter readings}
Globe at Night, run by NOIRLab, has collected observations of the night sky from volunteers since 2006 and publishes them under CC BY 4.0. Some observers report a reading from a Sky Quality Meter, a hand-held photometer that gives the brightness of the sky overhead in magnitudes per square arcsecond. I kept readings between 15 and 22.4 magnitudes made under a sky reported as clear, between 8 pm and 4 am local time, with valid coordinates: {{meta.gan_sqm_all|int}} readings from 2006 to 2025, of which {{model.n_fit|int}} were made in 2019 to 2025 with the standard Globe at Night form.

\subsection{Places and boundaries}
Dark Sky Places are taken from the table in the English Wikipedia article ``Dark-sky preserve'' (CC BY-SA 4.0), placed with the coordinates of each place's own article ({{dsp_world.placed|int}} of {{dsp_world.n|int}}), plus Wikidata items of dark sky place classes not in the table. Towns and villages come from OS Open Names for Great Britain and GeoNames for Northern Ireland; cities from GeoNames. Country outlines are Natural Earth; UK counties, Scottish council areas and Northern Ireland districts come from Office for National Statistics boundaries.

\subsection{Data not used}
The atlas of Falchi et al. (2016) is distributed by GFZ Data Services under a non-commercial licence, and its download carries a note that further distribution is prohibited; I therefore use only the figures printed in the paper (its Tables 1 and 2) for comparison. I found no open licence for the Night Blight maps of England published by CPRE in 2016, and did not use them.

\section{Methods}

\subsection{A kernel model of sky glow}
For each lit cell $j$, upward light $E_j$ is its mean radiance times its lit area. The artificial brightness of the sky overhead at a place $x$ is
\[ B(x) = k \sum_j E_j \, (d_{xj}^2 + d_0^2)^{-\alpha/2}, \]
where $d_{xj}$ is the distance in km, $d_0 = 1$ km keeps the kernel finite next to a source, and the sum runs over all cells within {{model.r_max|int}} km. This is the power law of distance behind Walker's (1977) rule of thumb, not a radiative transfer model: it ignores terrain, altitude, aerosols, the angular emission of lamps and their colour. The convolution is computed with fast Fourier transforms in latitude bands of two degrees, with the east to west cell width of the band's central latitude. On the 2.5 arc minute grid the kernel between two cells is the mean over $5 \times 5$ sub-points of each, so that a cell's value is its average brightness. For the UK grid, light from cells outside the United Kingdom, Ireland and the Isle of Man is convolved on the world grid and added.

\subsection{Calibration}
The exponent $\alpha$ (tried at 2.0, 2.3, 2.6 and 2.9), an optional exponential fall-off (none, 50 or 100 km) and the scale $k$ were fitted to the Globe at Night readings of 2019 to 2025. Model brightness is turned into a total sky brightness by adding the natural sky of 174 $\mu$cd/m\textsuperscript{2} (22.0 mag/arcsec\textsuperscript{2}, as in Falchi et al. 2016). I fitted $k$ by least absolute error in magnitudes, which resists the outliers that hand-held readings contain (thin cloud, moonlight, a nearby lamp), and only to readings darker than {{model.fit_min|f1}} magnitudes, because cells of 4.6 km cannot represent the brightest city-centre streets. The best form was chosen by five-fold spatial cross-validation in blocks of ten degrees.

\subsection{Thresholds and exposure}
Sky classes follow Falchi et al. (2016): artificial brightness of at most 1.74 $\mu$cd/m\textsuperscript{2} (1\% of the natural sky) is pristine; above 13.9 (8\%) the sky is light polluted; above 87 (half the natural sky); above 688 the Milky Way is hidden; and above 3,000 the eye cannot fully adapt to the dark. People and land are counted in each class, on the 2.5 arc minute grid for countries and on the 30 arc second grid for the United Kingdom.

\subsection{Distances and change}
For every city of a million or more people (GeoNames points, merging places within 25 km of a larger one), I found the nearest land cell below the Milky Way threshold, and the nearest below the light pollution threshold. For UK cities the search uses the 1 km grid. To describe change, the UK model was run with the lights of 2015 and of 2023 against the same 2023 population.

\section{Results}

\subsection{How well the model fits}
The best form has $\alpha$ = {{model.alpha|f1}} and no extra fall-off with distance. In cross-validation it misses readings darker than {{model.fit_min|f1}} magnitudes by {{model.cv18_mae|f2}} magnitudes on average (mean bias {{model.cv18_bias|f2}}; root mean square {{model.cv18_rmse|f2}}; correlation {{model.cv18_r|f2}}; Figure \ref{fig:fit}). Over all readings, including the bright ones, the root mean square error is {{model.cv_rmse|f2}}: the model badly underestimates the brightest city centres. Readings from 2012 to 2018, which were not used in the fit, are matched about as well (correlation {{model.older.r|f2}}). On the {{model.uk_fine.n|int}} readings inside the United Kingdom, the 1 km model misses by {{model.uk_fine.mae|f2}} magnitudes (correlation {{model.uk_fine.r|f2}}).

\begin{figure}[tbp]\centering
\includegraphics[width=\textwidth]{fig_fit.png}
\caption{Model against Globe at Night Sky Quality Meter readings, 2019 to 2025. Left: density of readings. Right: median model error in half-magnitude bins of the measured value; the fit uses only readings darker than {{model.fit_min|f1}}.}\label{fig:fit}
\end{figure}

\subsection{The world}
By the model, {{world.pop_mw|pct}} of the world's people (about {{world.people_mw_bn|f1}} billion) live under skies that hide the Milky Way, {{world.pop_polluted|pct}} under light polluted skies and {{world.pop_no_dark|pct}} where the eye cannot fully adapt to the dark (Figure \ref{fig:world}). The Milky Way is hidden over {{world.land_mw|pct}} of the land between 75\textdegree\ N and 60\textdegree\ S, and {{world.land_polluted|pct}} of that land is light polluted. The atlas of Falchi et al. (2016) gave {{falchi.pop_mw|pct}}, {{falchi.pop_polluted|pct}} and {{falchi.pop_no_dark|pct}} of people and {{falchi.land_polluted|pct}} of land.

\begin{figure}[tbp]\centering
\includegraphics[width=\textwidth]{fig_world_map.png}
\caption{Modelled artificial zenith sky brightness, 2023, at 2.5 arc minutes. Colours run from navy (under 1\% of the natural sky) through violet and rose (the Milky Way is hidden beyond about four times the natural sky) to amber and white.}\label{fig:world}
\end{figure}

Country rankings agree closely with the atlas (Figure \ref{fig:countries}, Table \ref{tab:countries}): Spearman's $\rho$ = {{falchi.compare.spearman|f2}} across {{falchi.compare.n|int}} countries of a million or more people, with my shares higher by {{falchi.mean_diff_pts|int}} percentage points on average. In {{countries_mw_over_half|int}} of {{countries_n|int}} such countries most people cannot see the Milky Way from home, and in {{countries_mw_over_90|int}} more than nine in ten cannot, among them Hong Kong, Singapore, the Gulf states, Israel and Egypt. In Eritrea, the Central African Republic, South Sudan, Sierra Leone and Ethiopia nearly everyone can. Ukraine is one of the few large countries where the 2023 figure is well below the atlas ({{country.UKR.pop_mw|pct}} against {{country.UKR.falchi_pop_mw|pct}}).

{{TABLE_COUNTRIES}}

\begin{figure}[tbp]\centering
\includegraphics[width=0.86\textwidth]{fig_countries.png}
\caption{Countries of 30 million people or more: share of people under skies that hide the Milky Way. Bars: this model, 2023. Ticks: Falchi et al. (2016).}\label{fig:countries}
\end{figure}

For cities of a million or more, the nearest Milky Way sky is a median {{cities_1m_median_mw_km|int}} km from the centre ({{cities_1m_n|int}} cities). The largest sprawls need more: {{city.New York City.mw_km|int}} km from New York, {{city.Moscow.mw_km|int}} km from Moscow, {{city.Delhi.mw_km|int}} km from Delhi and {{city.Shanghai.mw_km|int}} km from Shanghai (Figure \ref{fig:cities}).

\subsection{The United Kingdom}
At 1 km, {{uk.pop_mw|pct}} of the United Kingdom's people live under a Milky Way hiding sky, about {{uk.people_mw_m|f1}} million, but only {{uk.land_mw|pct}} of its land does (Figure \ref{fig:uk}, Table \ref{tab:nations}). The same population on the 4.6 km world grid gives {{uk.coarse_pop_mw|pct}}, a measure of how much the coarse grid understates exposure. The atlas gave {{country.GBR.falchi_pop_mw|pct}} of people for 2014.

{{TABLE_NATIONS}}

\begin{figure}[tbp]\centering
\includegraphics[width=0.78\textwidth]{fig_uk_map.png}
\caption{Modelled artificial zenith sky brightness over Great Britain and Ireland, 2023, at 30 arc seconds.}\label{fig:uk}
\end{figure}

The darkest areas, by median sky over their land, are the Western Isles (Na h-Eileanan Siar, {{uk_area.Na h-Eileanan Siar.land_mag|f2}} mag/arcsec\textsuperscript{2}), Highland, Orkney, Shetland and Argyll and Bute; in England they are Northumberland ({{uk_area.Northumberland.land_mag|f2}}), Devon, Cumbria, Cornwall and Dorset. The brightest are Greater London ({{uk_area.Greater London.land_mag|f2}}), Glasgow City and the West Midlands. In {{uk_councils_all_hidden|int}} of {{uk_councils_n|int}} local authority areas fewer than one resident in a hundred can see the Milky Way from home. The Milky Way should be visible from the centre of {{uk_kind_mw_visible.City|pct}} of cities in OS Open Names, {{uk_kind_mw_visible.Town|pct}} of towns, {{uk_kind_mw_visible.Village|pct}} of villages and {{uk_kind_mw_visible.Hamlet|pct}} of hamlets. The town with the darkest sky is Hugh Town in the Isles of Scilly ({{uk_towns_darkest.0.mag|f2}}). From central London the nearest Milky Way sky is {{uk_cities.London.mw_km|int}} km away, near Toys Hill in Kent; from Birmingham {{uk_cities.Birmingham.mw_km|int}} km, Manchester {{uk_cities.Manchester.mw_km|int}} km and Glasgow {{uk_cities.Glasgow.mw_km|int}} km.

All {{dsp_uk_n|int}} UK Dark Sky Places that could be placed lie below the Milky Way threshold in the model, from the Isle of Coll ({{dsp_uk.0.mag|f2}}) to Moore's Reserve in the South Downs ({{dsp_uk.16.mag|f2}}). Worldwide, {{dsp_world.mw_visible|pct}} of placed Dark Sky Places do.

\begin{figure}[tbp]\centering
\includegraphics[width=\textwidth]{fig_cities.png}
\caption{Distance from the centre of the world's largest cities (left) and the UK's largest cities (right) to the nearest land where the model sky is below the Milky Way threshold.}\label{fig:cities}
\end{figure}

\subsection{Change in the UK, 2015 to 2023}
Total VIIRS radiance over Great Britain fell by {{change.gbr_pct|abs1}}\% between 2015 and 2023, while over Ireland it rose by {{change.irl_pct|f1}}\% (Figure \ref{fig:change}). The largest falls were in Tyne and Wear ({{uk_light_fall.0.pct|abs1}}\%), Inverclyde and Northumberland; radiance rose in {{uk_areas_light_up|int}} of {{uk_areas_n|int}} areas, most in the Western Isles and Cornwall. With 2015 lights the model puts {{change.pop_mw_2015|pct}} of UK residents under a Milky Way hiding sky, with 2023 lights {{change.pop_mw_2023|pct}}; {{change.land_dimmer|pct}} of land has a sky at least 5\% dimmer and {{change.land_brighter|pct}} at least 5\% brighter.

\begin{figure}[tbp]\centering
\includegraphics[width=\textwidth]{fig_change.png}
\caption{Left: total VIIRS radiance over Great Britain and Ireland, 2015 = 100. Right: modelled change in UK sky glow, 2023 lights against 2015.}\label{fig:change}
\end{figure}

\section{Discussion}
A model this simple reproduces most of the variation in measured sky brightness outside city centres, and its country ranking is close to that of a far more elaborate atlas. Its level is higher. Several differences pull in that direction: lights from 2023 rather than 2014; a population grid constrained to built-up areas, which places people under their town's lights; and calibration readings that volunteers mostly take in the evening, when more lights are on than at the satellite's overpass. Others pull the other way: on a 4.6 km grid people share their cell's average sky, which understates exposure in towns, as the UK comparison shows. Without the atlas grid itself, which cannot be redistributed, I cannot separate these effects. The honest summary is a range: between about a third and a half of humanity cannot see the Milky Way from home.

The fall in the UK's radiance seen from space is consistent with the wholesale replacement of orange sodium street lamps with white LEDs, which are often dimmer and better shielded. But VIIRS sees little light below 500 nm, where white LEDs emit much of their energy and where light scatters most strongly in the air, so a street can look darker from space and brighter to the eye. Volunteers' counts of visible stars point the other way: Kyba et al. (2023) found from Globe at Night reports that the visibility of stars fell rapidly between 2011 and 2022, equivalent to a sky brightening of about 10\% a year. The satellite record should therefore not be read as evidence that British skies have become darker.

\section{Limitations}
The kernel ignores terrain shading, altitude, air clarity, the angle and colour of emission, snow and cloud. The calibration readings are concentrated in North America and Europe and are noisy; distinctions near natural darkness (the 1\% and 8\% thresholds) are within that noise, so the Milky Way threshold carries most weight here. The model underestimates the brightest city centres; I report them as ``never truly dark'' rather than with precise values. Kosovo has no light file of its own in the WorldPop set and is excluded from country results. Dark Sky Places are points at the centres of large areas.

\section{Conclusion}
Open satellite lights, an open population grid, citizen sky meter readings and a two-parameter model are enough to map where the Milky Way can still be seen, to within about half a magnitude outside city centres. By this model about half of the world's people and {{uk.pop_mw|pct}} of the United Kingdom's cannot see it from home, though in the UK it remains visible over {{uk.land_mw_visible|pct}} of the land, and from most large British cities a sky that shows it lies within an hour's drive.

\section*{Data and code availability}
Code, derived data and the interactive page: \url{https://github.com/Tristan-Hotham/dark-skies} and \url{https://maps.tristanhotham.com/dark-skies/}. Archived release with this paper: \url{https://doi.org/10.5281/zenodo.23757708}. Derived data and text are CC BY 4.0 (the Dark Sky Places table CC BY-SA 4.0); code is MIT.

\section*{AI statement}
I wrote the text, with drafting and editing assistance from Claude (Anthropic), which also helped prepare the analysis code. I reviewed and checked all text and results.

\section*{References}
\small
\begin{description}\setlength{\itemsep}{2pt}
\item Bondarenko, M., Priyatikanto, R., Tejedor-Garavito, N. et al. (2025) The spatial distribution of population in 2015 to 2030 at a resolution of 30 arc seconds, R2025A v1. WorldPop, University of Southampton. \url{https://doi.org/10.5258/SOTON/WP00845}
\item Cinzano, P. and Falchi, F. (2012) The propagation of light pollution in the atmosphere. Monthly Notices of the Royal Astronomical Society 427: 3337 to 3357.
\item Elvidge, C.D., Zhizhin, M., Ghosh, T., Hsu, F.-C. and Taneja, J. (2021) Annual time series of global VIIRS nighttime lights derived from monthly averages: 2012 to 2019. Remote Sensing 13(5): 922. \url{https://doi.org/10.3390/rs13050922}
\item Falchi, F., Cinzano, P., Duriscoe, D., Kyba, C.C.M., Elvidge, C.D., Baugh, K., Portnov, B.A., Rybnikova, N.A. and Furgoni, R. (2016) The new world atlas of artificial night sky brightness. Science Advances 2(6): e1600377. \url{https://doi.org/10.1126/sciadv.1600377}
\item Globe at Night (2006 to 2025) Campaign data. NOIRLab. \url{https://globeatnight.org/maps-data/} (CC BY 4.0).
\item Kyba, C.C.M., Alt\i nta\c{s}, Y.\"O., Walker, C.E. and Newhouse, M. (2023) Citizen scientists report global rapid reductions in the visibility of stars from 2011 to 2022. Science 379: 265 to 268.
\item Walker, M.F. (1977) The effects of urban lighting on the brightness of the night sky. Publications of the Astronomical Society of the Pacific 89: 405 to 409.
\item Woods, D., McKeen, T., Cunningham, A., Priyatikanto, R., Sorichetta, A., Tatem, A.J. and Bondarenko, M. (2024) WorldPop high resolution, harmonised annual global geospatial covariates, version 1.0. WorldPop, University of Southampton. \url{https://doi.org/10.5258/SOTON/WP00772}
\end{description}
\end{document}
